% ===================================================================
% SKILL: İÇERİK KUTULARI (BOXES) STANDARDI
% ===================================================================
% TANIM:
% Sunumda içeriğin türüne göre (Matematik, Uyarı, Örnek, Kod) doğru 
% kutu (block) yapısını kullanmak, öğrencinin algısını yönetir.
% 
% KURALLAR:
% 1. Matematik/Tanım -> Standart Mavi Blok (block)
% 2. Kritik/Uyarı    -> Kırmızı Blok (alertblock)
% 3. Örnek/Senaryo   -> Yeşil Blok (exampleblock)
% 4. Kod/Uzun Metin  -> Font küçültülmüş Blok (Taşmayı önlemek için)
% ===================================================================

\begin{frame}{İçerik Kutuları Örnekleri}

    % 1. MATEMATİK VE TEORİ (STANDART MAVİ BLOK)
    % Tanımlar, teoremler, formüller ve temel kavramlar için.
    \begin{block}{Tanım: Olasılık Yoğunluk Fonksiyonu}
        Bir sürekli rassal değişkenin belirli bir aralıkta değer alma olasılığını veren fonksiyondur.
        \[ P(a \le X \le b) = \int_{a}^{b} f(x) \,dx \]
    \end{block}

    \vspace{1em}

    % 2. KRİTİK UYARI (KIRMIZI ALERT BLOK)
    % Hatalar, dikkat edilmesi gerekenler, istisnalar için.
    \begin{alertblock}{Kritik Hata: Korelasyon vs Nedensellik}
        Korelasyonun yüksek olması ($r \approx 1$), bir değişkenin diğerinin \textbf{sebebi} olduğu anlamına gelmez! Bu ``Spurious Correlation'' tuzağıdır.
    \end{alertblock}

    \vspace{1em}

    % 3. ÖRNEK VE UYGULAMA (YEŞİL EXAMPLE BLOK)
    % Senaryolar, vaka analizleri ve pratik örnekler için.
    \begin{exampleblock}{Örnek Senaryo: Üretim Hattı}
        Bir fabrikada üretilen vidaların \%2'si kusurludur. Rastgele seçilen 100 vidadan en az 3 tanesinin kusurlu olma olasılığı nedir?
    \end{exampleblock}

\end{frame}

% 4. KOD BLOĞU (TAŞMA KORUMALI)
% Kodlar genellikle uzundur ve taşar. Mutlaka font küçültülmelidir.
% [fragile] parametresi kod bloğu içeren frame için ZORUNLUDUR.
\begin{frame}[fragile]{Kod Bloğu Örneği}
    
    \begin{block}{Python: Veri Analizi Kodu}
        % Font boyutu \small, \footnotesize veya \scriptsize olarak ayarlanır.
        % Bu sayede kodlar kutunun dışına taşmaz.
        \footnotesize 
        \begin{verbatim}
import pandas as pd
import numpy as np

# Veri setini yükle ve özetle
df = pd.read_csv('data.csv')

# Eksik verileri kontrol et (Kritik Adım)
if df.isnull().sum() > 0:
    print("Uyarı: Eksik veri bulundu!")
    df.dropna(inplace=True)

print(df.describe())
        \end{verbatim}
    \end{block}
    
    \vspace{0.5em}
    \centering
    \textit{\scriptsize Not: Kod bloklarında fontu küçültmek (footnotesize) taşmayı önler.}
\end{frame}
